% Generated draft: requires mathematical and source review.
\documentclass[11pt]{article}
\usepackage[margin=25mm]{geometry}
\usepackage{fontspec}
\setmainfont{Latin Modern Roman}
\usepackage{amsmath,amssymb,mathtools,unicode-math}
\setmathfont{Latin Modern Math}
\usepackage{enumitem,needspace,xurl,hyperref}
\setlength{\parindent}{0pt}
\setlength{\parskip}{5pt}
\setlength{\emergencystretch}{3em}
\newcommand{\N}{\mathbb{N}}
\newcommand{\Z}{\mathbb{Z}}
\newcommand{\Q}{\mathbb{Q}}
\newcommand{\R}{\mathbb{R}}
\newcommand{\C}{\mathbb{C}}
\newcommand{\F}{\mathbb{F}}
\newcommand{\A}{\mathbb{A}}
\newcommand{\PP}{\mathbb{P}}
\newcommand{\E}{\mathbb{E}}
\newcommand{\eps}{\varepsilon}
\newcommand{\dd}{\,\mathrm{d}}
\begin{document}
\section*{UnsolvedMath ID: 801142}
\textbf{Source problem code:} catalogue entry
\subsection{Definitions and mathematical statement}
A type \(\mathrm{II}_1\) factor is a von Neumann algebra \(M\) with center \(Z(M)=\C 1\), no nonzero minimal projections, and a faithful normal tracial state \(\tau_M\) satisfying \(\tau_M(1)=1\). In a common continuous first-order language \(L\) for tracial von Neumann algebras, write \(\sigma^M\) for the value of an \(L\)-sentence \(\sigma\) in \(M\). Elementary equivalence means
\[
M\equiv_L N
\quad\Longleftrightarrow\quad
\sigma^M=\sigma^N\text{ for every }L\text{-sentence }\sigma.
\]
The source asks whether two type \(\mathrm{II}_1\) factors can have the same ``2-quantifier theory'' without being elementarily equivalent. The supplied record does not define the language or the intended quantifier-complexity convention. Accordingly, let \(\mathcal S_2\) denote the class of sentences intended by that phrase, pending clarification. Relative to a specified \(L\) and \(\mathcal S_2\), the requested existence statement is
\[
\begin{gathered}
\exists\, M,N\text{ of type }\mathrm{II}_1:\\
\bigl(\forall\sigma\in\mathcal S_2,\ \sigma^M=\sigma^N\bigr)
\quad\text{and}\quad
\bigl(\exists\text{ an }L\text{-sentence }\varphi,\ \varphi^M\ne\varphi^N\bigr).
\end{gathered}
\]
This is a formulation schema until the meaning of ``2-quantifier theory'' is fixed: the record does not say whether it counts individual quantifiers, quantifier blocks, quantifier depth, or alternations, nor whether it includes both quantifier orders. No separability assumption or requirement that \(M\) and \(N\) themselves be ultraproducts is stated. The supplied source points to Proposition 7 in [OzawaSolidActa] and also to [HURigid], and describes the problem as still open; these are supplied source qualifications, not independently verified conclusions.
\subsection{Short English statement}
Do there exist two type \(\mathrm{II}_1\) factors that agree on every sentence in the source's intended 2-quantifier fragment but disagree on some first-order sentence? The fragment and language require clarification before this becomes a uniquely specified mathematical question. The supplied record labels the problem open.
\subsection{Sources}
\begin{itemize}
\item Supplied local UnsolvedMath registry record, permanent ID \(801142\), exact source problem code \texttt{catalogue entry}, titled ``Ultraproduct \(\mathrm{II}_1\) Factors --- Problem 44.'' Attribution: Elayavalli, proposed year 2025; registry status: open.
\item Supplied source-list attribution: Elayavalli, \emph{50 Open Problems: Ultraproduct \(\mathrm{II}_1\) Factors} (2025), Problem 44. Supplied URL: \url{https://arxiv.org/abs/2511.20377}. The record reports access on 2026-10-16 and extraction by \texttt{arxiv\_latex}; these are supplied provenance details, not an independent source visit or verification.
\item The supplied statement directs the reader to Proposition 7 in [OzawaSolidActa], with an additional reference to [HURigid]. No bibliographic entries or links resolving these two reference labels are supplied, so their identities, hypotheses, and relevance remain unverified.
\end{itemize}
\end{document}
